% radiosynth - From MHD to MHz: thermal gyroresonance radio emission, step-by-step calculation
% Copyright (C) 2026 Soumitra Hazra
% Compile with: pdflatex radiosynth_gyroresonance_steps.tex (twice)
\documentclass[11pt,a4paper]{article}
\usepackage[T1]{fontenc}
\usepackage{lmodern}
\usepackage{helvet}
\renewcommand{\familydefault}{\sfdefault}
\usepackage[margin=2.4cm,bottom=2.6cm]{geometry}
\usepackage{amsmath,amssymb,bm}
\usepackage{booktabs}
\usepackage{xcolor}
\usepackage{enumitem}
\usepackage{titlesec}
\usepackage{fancyhdr}
\usepackage{needspace}
\usepackage[colorlinks=true,linkcolor=blue!50!black,urlcolor=blue!50!black]{hyperref}

\definecolor{head}{RGB}{31,73,125}
\definecolor{boxbg}{RGB}{238,244,250}
\titleformat{\section}{\needspace{6\baselineskip}\color{head}\large\bfseries}{\thesection.}{0.4em}{}
\titleformat{\subsection}{\color{head}\normalsize\bfseries\itshape}{\thesubsection}{0.5em}{}
\titlespacing*{\section}{0pt}{14pt}{5pt}
\titlespacing*{\subsection}{0pt}{9pt}{3pt}
\pagestyle{fancy}\fancyhf{}\renewcommand{\headrulewidth}{0pt}
\fancyfoot[R]{\small\color{gray}Page \thepage}
\setlength{\parskip}{5pt}\setlength{\parindent}{0pt}
\setlength{\emergencystretch}{2.5em}
\setlist{itemsep=1pt,topsep=3pt}

\newcommand{\vect}[1]{\bm{#1}}
\newcommand{\dd}{\mathrm{d}}
\newcommand{\kB}{k_\mathrm{B}}
\newcommand{\nuB}{\nu_B}
\newcommand{\code}[1]{\texttt{#1}}
\newcommand{\example}[1]{\par\noindent\colorbox{boxbg}{\parbox{\dimexpr\linewidth-2\fboxsep}{\small\textbf{Numerical example.} #1}}\par}
\newcommand{\whatitmeans}[1]{\par{\color{head}\textbf{What this means:}} #1\par}
\newcommand{\incode}[1]{\par{\small\color{head!80!black}\textbf{In radiosynth:} #1}\par}

\hypersetup{pdftitle={Thermal Gyroresonance Radio Emission in radiosynth},pdfauthor={Soumitra Hazra}}

\begin{document}

\setcounter{section}{-1}
\begin{center}
{\LARGE\bfseries\color{head} Thermal Gyroresonance Radio Emission in radiosynth\par}
\end{center}
\vspace{2pt}
\textbf{Step-by-step calculation from a 3D MHD simulation}\\
With separate derivation subsections for the Appleton--Hartree formula and the harmonic-layer
optical depth.

Soumitra Hazra --- \textsf{radiosynth} 2.2.0, \emph{From MHD to MHz}, November 2026 ---
\url{https://github.com/soumitrahazra/radiosynth}

The equations are written in Gaussian-cgs units; for SI, replace $e^2$ by $e^2/(4\pi\varepsilon_0)$
(radiosynth itself computes in SI). A single worked example runs through all steps: a ray at
5\,GHz crossing the third-harmonic layer in a corona with $n_e=10^9$\,cm$^{-3}$ and
$T_e=2$\,MK, at $\theta=60^\circ$ to the field. Every number in it was computed with radiosynth.

\textbf{How to read this note.} Each step has the same shape: a short explanation, the equations
(with every intermediate line of the derivations in Steps~4 and 5), a sentence on \emph{what this means},
the worked numbers, and the function of radiosynth that does it. You can follow the physics
without the code, or the code without the derivations.

%======================================================================
\section{Start from the MHD model and set up the rays}
%======================================================================
The model supplies, at any point, the electron density $n_e$, the electron temperature $T_e$ and
the magnetic field $\vect B$ (in gauss). For each image pixel, radiosynth uses one
\emph{straight} line of sight: refraction is negligible where gyroresonance emission forms,
because there the observing frequency is far above the plasma frequency. Let $\hat{\vect k}$ be
the unit vector along the line, \emph{towards the observer}. The line is sampled at nodes
$k=0,1,\ldots,N$ spaced by $\Delta\ell$, with node 0 nearest to the observer, from where the line
enters the model to where it leaves it or reaches the photosphere. At every node:
\[
  B_k=|\vect B_k|,\qquad \cos\theta_k=\frac{\vect B_k\cdot\hat{\vect k}}{B_k},\qquad
  n_{e,k},\quad T_{e,k}.
\]
By default $\Delta\ell$ is half the smallest cell of the model grid. It must resolve the variation
of the model, not the thin layers, which are found between the nodes (Step~2).
\incode{\code{common/straight.py}: \code{segment\_ends}, \code{default\_step}, \code{iter\_samples}.}

%======================================================================
\section{Choose an observing frequency and detect the resonance layers}
%======================================================================
For thermal electrons, the nonrelativistic gyrofrequency is
\begin{equation}
  \nuB=\frac{eB}{2\pi m_ec}\simeq2.80\times10^6\,B[\text{G}]\ \text{Hz}.
\end{equation}
At harmonic $s$, the resonance condition and the field strength it requires are
\begin{equation}
  \nu=s\,\nuB,\qquad B_s=\frac{\nu}{2.80\times10^6\,s}\ \text{G}.
\end{equation}
\whatitmeans{an electron spirals around the field line $\nuB$ times per second, and it can absorb or emit a
wave whose frequency is $\nuB$ or an integer multiple of it. Because $B$ falls with height above a sunspot,
each harmonic $s$ picks out one surface of constant $|B|$.}
At 5\,GHz: $s=1$: 1786\,G; $s=2$: 893\,G; $s=3$: 595\,G; $s=4$: 447\,G. A ray can cross several
of these surfaces of constant $|B|$.

Between nodes $k$ and $k+1$ a layer of harmonic $s$ is present when
\begin{equation}
  (s\,\nu_{B,k}-\nu)(s\,\nu_{B,k+1}-\nu)<0
  \quad\Longleftrightarrow\quad (B_k-B_s)(B_{k+1}-B_s)<0 .
\end{equation}
All harmonics $s=1\ldots s_\text{max}$ (default 10) are tested in every step, and several layers can
fall into one step. This assumes $B$ varies smoothly between nodes; use smaller steps if a
harmonic surface could be crossed twice in one step.

%======================================================================
\section{Locate each crossing and interpolate the plasma}
%======================================================================
For a crossing between $B_k$ and $B_{k+1}$, radiosynth interpolates linearly in $\nuB$ (that is, in
$B$):
\begin{equation}
  u=\frac{\nu/s-\nu_{B,k}}{\nu_{B,k+1}-\nu_{B,k}}=\frac{B_s-B_k}{B_{k+1}-B_k},\qquad
  \ell_c=\ell_k+u\,\Delta\ell .
\end{equation}
The plasma parameters on the layer are interpolated with the same weight,
\begin{equation}
  q_c=q_k+u\,(q_{k+1}-q_k),\qquad q\in\{n_e,\ T_e,\ \cos\theta,\ X\},
\end{equation}
where $X=(\nu_p/\nu)^2$ is defined in Step~4. Note that radiosynth interpolates $\cos\theta$ between
the nodes, rather than recomputing it from an interpolated field vector; for steps that resolve
the field the two agree.

\example{$B_k=500$\,G, $B_{k+1}=650$\,G and $B_3=595.4$\,G give $u=95.4/150=0.636$. For a
1000\,km step, the crossing is 636\,km from node $k$.}

%======================================================================
\section{Calculate the magnetic scale length}
%======================================================================
The integrated opacity depends on how fast the field strength changes along the ray:
\begin{equation}
  L_B=\left|\frac{B}{\dd B/\dd\ell}\right|_c\simeq\frac{B_s\,\Delta\ell}{|B_{k+1}-B_k|}.
\end{equation}
\whatitmeans{a wave crossing the layer stays in resonance only while $s\nuB$ is within the thin line width of
$\nu$. If the field changes slowly (large $L_B$), that takes a longer path, so more electrons absorb: the
optical depth is proportional to $L_B$ (Step~5).}
\example{$B_3=595.4$\,G and a 150\,G change across 1000\,km give $L_B=3969$\,km.}
\incode{\code{gyroresonance/transfer.py}, function \code{transfer\_chunk}: the test \code{cross},
the weight \code{u}, \code{lerp}, and \code{l\_b}.}

%======================================================================
\section{Derivation subsection A: the Appleton--Hartree magneto-ionic formula}
%======================================================================
\subsection{Definitions and physical assumptions}
In a cold, collisionless, magnetised electron plasma (ions at rest), linearising the electron
equation of motion and Maxwell's equations gives a dielectric tensor. Solving the dispersion
relation for transverse waves yields two eigenmodes, the ordinary (o) and extraordinary (x)
modes. Define
\begin{equation}
  X=\frac{\nu_p^2}{\nu^2},\qquad Y=\frac{\nuB}{\nu},\qquad
  \nu_p^2=\frac{n_ee^2}{\pi m_e}\quad\Big(\text{SI: } \frac{n_ee^2}{4\pi^2\varepsilon_0m_e}\Big).
\end{equation}
The refractive index is $n=ck/\omega$, and $\theta$ is the angle between the wave vector and the
magnetic field.

\subsection{Algebraic route from the dielectric tensor}
In a frame with $\vect B$ along $z$, the cold-plasma dielectric tensor has the Stix components
\begin{equation}
  R=1-\frac{X}{1-Y},\quad L=1-\frac{X}{1+Y},\quad S=\frac{R+L}{2},\quad D=\frac{R-L}{2},\quad P=1-X,
\end{equation}
with diagonal elements $(S,S,P)$, $\epsilon_{xy}=-iD$, $\epsilon_{yx}=iD$, and all other entries zero.
For a plane wave $\propto e^{i(\vect k\cdot\vect r-\omega t)}$, Maxwell's equations give
$\vect n\times(\vect n\times\vect E)+\boldsymbol\epsilon\cdot\vect E=0$ with $\vect n=c\vect k/\omega$. Taking $\vect n=n(\sin\theta,0,\cos\theta)$ in the $x$--$z$ plane, this is
$\mathbf M\,\vect E=0$ with the wave matrix
\begin{equation}
  \mathbf M=\begin{pmatrix}
    S-n^2\cos^2\theta & -iD & n^2\sin\theta\cos\theta\\
    iD & S-n^2 & 0\\
    n^2\sin\theta\cos\theta & 0 & P-n^2\sin^2\theta
  \end{pmatrix}.
\end{equation}
A wave exists only if $\det\mathbf M=0$. Expanding along the middle row and using $S^2-D^2=RL$,
\begin{align}
  \det\mathbf M&=(S-n^2)\Big[(S-n^2\cos^2\theta)(P-n^2\sin^2\theta)-n^4\sin^2\theta\cos^2\theta\Big]
              -D^2\,(P-n^2\sin^2\theta)\nonumber\\
  &=A\,n^4-B_\text{S}\,n^2+C=0,
\end{align}
which is a quadratic in $n^2$,
\begin{equation}
  A\,n^4-B_\text{S}\,n^2+C=0,\quad
  A=S\sin^2\theta+P\cos^2\theta,\quad
  B_\text{S}=RL\sin^2\theta+PS(1+\cos^2\theta),\quad
  C=PRL,
\end{equation}
with the two roots
\begin{equation}
  n^2=\frac{B_\text{S}\pm F}{2A},\qquad F^2=(RL-PS)^2\sin^4\theta+4P^2D^2\cos^2\theta .
\end{equation}
Now insert $R$, $L$, $S$, $D$, $P$ in terms of $X$ and $Y$. Three identities do most of the work:
\begin{equation}
  RL=\frac{(1-X)^2-Y^2}{1-Y^2},\qquad RL-PS=-\frac{XY^2}{1-Y^2},\qquad D=-\frac{XY}{1-Y^2},
\end{equation}
so that
\begin{equation}
  F=\frac{XY}{1-Y^2}\sqrt{Y^2\sin^4\theta+4(1-X)^2\cos^2\theta}.
\end{equation}
Dividing $B_\text{S}\pm F$ by $2A$ and simplifying, the two roots take the collisionless Appleton--Hartree form
\begin{equation}
  n_\pm^2=1-\frac{X(1-X)}{(1-X)-\tfrac12Y^2\sin^2\theta\pm\sqrt{\tfrac14Y^4\sin^4\theta+(1-X)^2Y^2\cos^2\theta}} .
  \label{eq:ah}
\end{equation}
\whatitmeans{in a magnetised plasma a radio wave can only travel in two ways, and each has its own
refractive index. The formula tells us which way is which, and where each one stops.}
(the root $(B_\text{S}+F)/2A$ is the upper sign). The branches are identified from their limits. Along the field ($\theta=0$) the denominator is
$(1-X)(1\pm Y)$, so $n_+^2=1-X/(1+Y)=L$ and $n_-^2=1-X/(1-Y)=R$. Across the field ($\theta=90^\circ$),
$n_+^2=1-X$. The upper sign is therefore the o-mode and the lower sign the x-mode, and for $Y<1$
the cutoffs ($n^2=0$) are
\begin{equation}
  \text{o-mode: } X\ge1,\qquad\qquad \text{x-mode: } X\ge1-Y .
\end{equation}
A non-propagating mode cannot transmit radiation from behind its cutoff. On the fundamental
layer ($Y=1$) the x-mode is cut off for any $X>0$.

\subsection{Polarisation of the modes}
The same eigenvectors give the polarisation. Neglecting the small longitudinal part, the electric
vector of mode $\sigma$ traces an ellipse with axial ratio $T_\sigma$, $T_oT_x=-1$, and
\begin{equation}
  |T_x|=\frac{\sqrt{Y^2\sin^4\theta+4(1-X)^2\cos^2\theta}-Y\sin^2\theta}{2(1-X)|\cos\theta|}
       =\frac{2(1-X)|\cos\theta|}{Y\sin^2\theta+\sqrt{Y^2\sin^4\theta+4(1-X)^2\cos^2\theta}}
\end{equation}
(the second form, used in the code, avoids the cancellation at small $\cos\theta$). Along the field
$|T_x|=1$: both modes are circular. Across the field $T_x=0$: the x-mode is linear with
$\vect E\perp\vect B$. The x-mode rotates in the same sense as the gyrating electrons. Seen by an
observer towards whom $\vect B$ points ($\vect B\cdot\hat{\vect k}>0$), the electrons, and hence the x-mode, rotate
counter-clockwise: right-hand circular polarisation in the IAU/IEEE convention, Stokes $V>0$.
\example{At the $s=3$ layer at 5\,GHz, $X=0.0032$ and $Y=1/3$, so $n^2_o=0.997$ and
$n^2_x=0.996$ at $60^\circ$: the refractive index is 1 to better than 0.5\,\%, which justifies setting
$n_\sigma=1$ in Step~5. The axial ratio is $|T_x|=0.78$.}

This is a cold-plasma propagation result. It does not by itself give the gyroresonance
absorption: that is supplied by the thermal electron distribution (Subsection~B).
\incode{\code{common/magnetoionic.py}: \code{magnetoionic\_params}, \code{refractive\_index2}
(Eq.~\ref{eq:ah}), \code{cutoff}, \code{axial\_ratio\_x}.}

%======================================================================
\section{Derivation subsection B: optical depth of a harmonic layer}
%======================================================================
\subsection{Thermal resonance width}
Thermal electrons move along the field, so the resonance is Doppler-shifted. For weakly
relativistic electrons,
\begin{equation}
  \nu\left(1-n_\sigma\beta_\parallel\cos\theta\right)\simeq\frac{s\,\nuB}{\gamma}.
\end{equation}
To leading order, with $n_\sigma\simeq1$, a normalised Gaussian represents the resonance:
\begin{equation}
  \phi_s(\nu)=\frac{1}{\sqrt\pi\,\Delta\nu_D}\exp\!\left[-\frac{(\nu-s\nuB)^2}{\Delta\nu_D^2}\right],
  \qquad
  \Delta\nu_D=\sqrt2\,\nu\,\beta_T|\cos\theta|,\qquad \beta_T^2=\frac{\kB T_e}{m_ec^2}.
\end{equation}
Near $\theta=90^\circ$ the first-order Doppler width vanishes and relativistic broadening takes over;
as shown in 5.4, the final optical depth does not depend on the width, so this does not change
the result.

\subsection{Low-temperature harmonic strength}
The cyclotron coupling involves the Bessel functions $J_{s-1}(s\beta_\perp\sin\theta)$. For small
arguments, $J_{s-1}(x)\simeq(x/2)^{s-1}/(s-1)!$. A Maxwellian distribution of perpendicular
velocities is
\begin{equation}
  f(\beta_\perp)\,\dd\beta_\perp\propto\exp\!\left[-\frac{\beta_\perp^2}{2\beta_T^2}\right]\beta_\perp\,\dd\beta_\perp .
\end{equation}
Its moments follow with the substitution $t=\beta_\perp^2/(2\beta_T^2)$, $\beta_\perp\,\dd\beta_\perp=\beta_T^2\,\dd t$:
\begin{equation}
  \left\langle\beta_\perp^{2m}\right\rangle
  =\frac{\int_0^\infty\beta_\perp^{2m}e^{-\beta_\perp^2/2\beta_T^2}\beta_\perp\,\dd\beta_\perp}{\int_0^\infty e^{-\beta_\perp^2/2\beta_T^2}\beta_\perp\,\dd\beta_\perp}
  =\frac{(2\beta_T^2)^m\int_0^\infty t^me^{-t}\,\dd t}{\int_0^\infty e^{-t}\,\dd t}=m!\,(2\beta_T^2)^m .
\end{equation}
Squaring the small-argument form, $J_{s-1}^2\simeq(s\sin\theta/2)^{2(s-1)}\beta_\perp^{2(s-1)}/[(s-1)!]^2$, and averaging with $m=s-1$:
\begin{equation}
  s\left\langle J_{s-1}^2\right\rangle\simeq s\,\frac{(s\sin\theta/2)^{2(s-1)}\,(s-1)!\,(2\beta_T^2)^{s-1}}{[(s-1)!]^2}
  =\frac{s^2}{s!}\left(\frac{s^2\beta_T^2\sin^2\theta}{2}\right)^{s-1}\left[1+O(s^2\beta_T^2)\right].
\end{equation}
\whatitmeans{the bracket is tiny ($\sim10^{-3}$) and is raised to the power $s-1$, so each higher harmonic is roughly a
thousand times weaker; and $\sin^2\theta$ means electrons hardly absorb a wave travelling along the field.}
The polarisation of the mode introduces the factor
\begin{equation}
  F_\sigma(\theta)=\frac{(1+T_\sigma\cos\theta)^2}{1+T_\sigma^2},
\end{equation}
with $T_\sigma$ from 4.3 evaluated on the layer ($Y=1/s$) and signed so that $T_x\cos\theta>0$. Along the field
$F_x=2$ and $F_o=0$; across it $F_x=1$ and $F_o=0$ at this order.
\example{For $s=3$ at $60^\circ$: $F_x=1.201$ and $F_o=0.049$. With $\beta_T^2=3.37\times10^{-4}$
(2\,MK), the factor $(s^2\beta_T^2\sin^2\theta/2)$ is $1.14\times10^{-3}$: each extra harmonic costs about
three orders of magnitude.}

\subsection{Local absorption coefficient}
The local absorption coefficient of mode $\sigma$ near harmonic $s$ is a line of fixed area
$A_{s\sigma}$ times the normalised profile:
\begin{equation}
  \kappa_{s\sigma}(\nu,\ell)=A_{s\sigma}(\ell)\,\phi_s\big[\nu-s\nuB(\ell)\big],\qquad
  A_{s\sigma}=\frac{\pi e^2}{m_ec}\,n_e\,\frac{s^2}{s!}\left(\frac{s^2\beta_T^2\sin^2\theta}{2}\right)^{s-1}F_\sigma(\theta),
\end{equation}
neglecting higher-order thermal terms and with $n_\sigma\simeq1$ (the full expression contains a factor
$n_\sigma^{2s-3}$). Numerically, $\pi e^2/(m_ec)=0.02654$\,cm$^2$\,s$^{-1}$.

\subsection{Integrate through the narrow layer}
At fixed $\nu$, the ray passes through a region where $s\nuB(\ell)$ sweeps across $\nu$. This region has
thickness $\delta\ell\simeq\sqrt2\beta_T|\cos\theta|\,L_B$, a few per cent of $L_B$, so $A_{s\sigma}$ and the field
gradient are nearly constant across it. Then
\begin{equation}
  \tau_{s\sigma}=\int\kappa_{s\sigma}\,\dd\ell=A_{s\sigma}\int\phi_s\big(\nu-s\nuB(\ell)\big)\,\dd\ell .
\end{equation}
Change the integration variable from the path $\ell$ to the local resonance frequency
$q=s\nuB(\ell)$. Then $\dd q=s\,(\dd\nuB/\dd\ell)\,\dd\ell$, i.e.\ $\dd\ell=\dd q/(s|\dd\nuB/\dd\ell|)$, and
\begin{equation}
  \int\phi_s\big(\nu-s\nuB(\ell)\big)\,\dd\ell=\frac{1}{s|\dd\nuB/\dd\ell|}\int\phi_s(\nu-q)\,\dd q=\frac{1}{s|\dd\nuB/\dd\ell|},
\end{equation}
because the profile has unit area. Hence
\begin{equation}
  \tau_{s\sigma}=\frac{A_{s\sigma}}{s\,|\dd\nuB/\dd\ell|}.
\end{equation}
On the resonance $s\nuB=\nu$ and $|\dd\nuB/\dd\ell|=\nuB/L_B$, so $s|\dd\nuB/\dd\ell|=\nu/L_B$. Substituting
gives the optical depth of the layer:
\begin{equation}
  \boxed{\;
  \tau_{s\sigma}=\frac{\pi e^2}{m_ec}\,\frac{n_eL_B}{\nu}\;\frac{s^2}{s!}
  \left(\frac{s^2\beta_T^2\sin^2\theta}{2}\right)^{s-1}F_\sigma(\theta)\;}
  \label{eq:tau}
\end{equation}
\whatitmeans{the layer is a thin wall whose opacity depends only on the plasma at the wall ($n_e$, $T_e$,
$\theta$), on how fast the field changes there ($L_B$), and on the mode. Its microscopic thickness drops out.}
The Gaussian width cancels because the whole resonance is integrated. This is why radiosynth can
treat each layer as a jump in optical depth instead of resolving its width.

\subsection{Interpretation and limits}
At fixed $s$ and $\nu$, $\tau_{s\sigma}\propto n_e\,L_B\,T_e^{s-1}\sin^{2s-2}\theta\,F_\sigma(\theta)$. Higher harmonics
are much thinner, and a layer seen along the field ($\sin\theta\to0$) is almost transparent. At the
fundamental the x-mode is cut off (4.2), so radiosynth sets $\tau_{1x}=0$ when $X>0$.
\example{Our layer ($n_e=10^9$\,cm$^{-3}$, $T_e=2$\,MK, $\theta=60^\circ$, $L_B=3969$\,km) has
$\tau_{3x}=4.92$ and $\tau_{3o}=0.200$. With $L_B=10^9$\,cm the same layer would have
$\tau_{3x}=12.4$ and $\tau_{3o}=0.50$; there the $s=2$ layer is far thicker
($\tau_{2x}\approx6600$, $\tau_{2o}\approx130$) and the $s=4$ layer thin ($\tau_{4x}\approx0.035$).}
\incode{\code{gyroresonance/opacity.py}: \code{mode\_factor} and \code{layer\_optical\_depth}
(Eq.~\ref{eq:tau}); \code{harmonic\_field} gives $B_s$.}

%======================================================================
\section{Add free--free absorption between layers}
%======================================================================
Between the layers, the plasma absorbs by free--free collisions. radiosynth uses the
Coulomb-logarithm form of Dulk (1985), evaluated at the mid-point of the step,
\begin{equation}
  \kappa_\text{ff}=9.78\times10^{-3}\,\frac{n_e^2}{\nu^2T_e^{3/2}}\big(24.5+\ln T_e-\ln\nu\big)\ \text{cm}^{-1}
  \qquad(T_e>2\times10^5\,\text{K}),
\end{equation}
split between the modes as
\begin{equation}
  \kappa_{\text{ff},x}=\frac{\kappa_\text{ff}}{(1-q)^2},\qquad
  \kappa_{\text{ff},o}=\frac{\kappa_\text{ff}}{(1+q)^2},\qquad q=\min(Y|\cos\theta|,\,0.9).
\end{equation}
Each step has $\Delta\tau_{\text{ff},\sigma}=\kappa_{\text{ff},\sigma}\Delta\ell$. Adding the layers crossed in the step,
\begin{equation}
  \Delta\tau_\sigma=\Delta\tau_{\text{ff},\sigma}+\sum_{s\,\in\,\text{crossings}}\tau_{s\sigma}.
\end{equation}
By Kirchhoff's law each mode emits $\nu^2\kB T_e/c^2$. When the parts of a step have different
temperatures, it glows at the optical-depth-weighted temperature
\begin{equation}
  T_\text{eff}=\frac{\Delta\tau_{\text{ff},\sigma}\,\bar T_e+\sum_s\tau_{s\sigma}T_{e,c}}{\Delta\tau_\sigma}.
\end{equation}
\example{For our step, $\Delta\tau_{\text{ff},x}=3.3\times10^{-4}$ and $\Delta\tau_{\text{ff},o}=1.7\times10^{-4}$ over
1000\,km: negligible next to the layer, so $\Delta\tau_x=4.92$, $\Delta\tau_o=0.200$ and $T_\text{eff}=2$\,MK.}
\incode{\code{bremsstrahlung/emission.py} (\code{gaunt="dulk"}) and
\code{bremsstrahlung/magnetized.py}; in \code{transfer\_chunk}, \code{dtau\_x}, \code{dtau\_o} and the sums
\code{st\_x}, \code{st\_o}.}

%======================================================================
\section{Solve the radiative transfer towards the observer}
%======================================================================
Along the ray, the intensity of one polarisation obeys $\dd I/\dd\tau=S-I$: absorption removes $I$ and the
plasma adds its source function $S=\nu^2\kB T_\text{eff}/c^2$. In brightness-temperature units
($T=Ic^2/\kB\nu^2$) this is $\dd T/\dd\tau=T_\text{eff}-T$. Multiplying by $e^{\tau}$ gives $\dd(Te^{\tau})/\dd\tau=T_\text{eff}e^{\tau}$, and
integrating across a step of optical depth $\Delta\tau$ with constant $T_\text{eff}$ gives the exact step formula.
For each polarisation, the transfer can therefore be done from the far end of the ray towards the
observer, one step at a time:
\begin{equation}
  T_{b,\text{out}}=T_{b,\text{in}}\,e^{-\Delta\tau}+T_\text{eff}\left(1-e^{-\Delta\tau}\right).
  \label{eq:rec}
\end{equation}
If $\Delta\tau\gg1$ the outgoing brightness approaches the local electron temperature; if $\Delta\tau\ll1$
and the background is negligible, the step adds about $T_\text{eff}\,\Delta\tau$.

radiosynth sums the same thing from the observer inwards. For two steps, with step 0 next to
the observer and nothing behind step 1, Eq.~\eqref{eq:rec} gives
\begin{equation}
  T=\underbrace{T_1\left(1-e^{-\Delta\tau_1}\right)}_{\text{leaves step 1}}e^{-\Delta\tau_0}+T_0\left(1-e^{-\Delta\tau_0}\right),
\end{equation}
so each step's emission is dimmed by all the steps in front of it. For any number of steps:
\begin{equation}
  T_p=\sum_{k=0}^{N-1}T_{\text{eff},k}\left(1-e^{-\Delta\tau_k}\right)\exp\!\Big(-\sum_{j<k}\Delta\tau_j\Big),
\end{equation}
where $T_p=I_pc^2/(\kB\nu^2)$ is the brightness temperature of one polarisation (each carries half
of the Rayleigh--Jeans intensity). The two forms are mathematically identical. A polarisation
stops at its first cutoff node (4.2): steps behind it contribute nothing. Optionally, lines that
end on the photosphere add a body temperature $T_\text{body}\,e^{-\tau_p}$.
\whatitmeans{an opaque step ($\Delta\tau\gg1$) shows its own temperature and hides everything behind it; a transparent one
adds a little of its temperature and lets the light from behind through.}
\example{With no emission behind the step, the x-mode gives $T=2\,\text{MK}\,(1-e^{-4.92})=1.99$\,MK
and the o-mode $T=2\,\text{MK}\,(1-e^{-0.200})=0.36$\,MK.}
\incode{\code{transfer\_chunk}: \code{alive} (cutoffs), \code{tau\_front}, \code{w = exp(-tau\_front) *
-expm1(-d)}, and the body term with the keyword \code{t\_body}.}

%======================================================================
\section{Combine the polarisations and form radio maps}
%======================================================================
Where $\vect B\cdot\hat{\vect k}$ changes sign along a ray (a quasi-transverse region), the radiation either keeps
its sense of circular polarisation (weak coupling, the default and the usual case at
microwaves) or keeps its mode identity (strong coupling). With weak coupling, radiosynth follows
right- ($R$) and left-hand ($L$) circular polarisation: in each step, $R$ uses the x-mode optical depth
where $\vect B\cdot\hat{\vect k}>0$ and the o-mode one elsewhere, and $L$ the opposite. With strong coupling it
follows the x- and o-modes and takes the sign of $V$ from the field direction at the observer's
end of the ray. Then
\begin{equation}
  I=I_R+I_L,\qquad V=I_R-I_L,\qquad P_c=V/I,
\end{equation}
and the Rayleigh--Jeans brightness temperatures are
\begin{equation}
  T_b=\frac{c^2I}{2\kB\nu^2}=\tfrac12\,(T_R+T_L),\qquad T_b(V)=\tfrac12\,(T_R-T_L).
\end{equation}
Repeating this for every pixel gives the brightness-temperature and polarisation maps, and the
flux density is $F=\sum_\text{pixels}I\,\Delta\Omega$ with $\Delta\Omega$ the solid angle of a pixel.
\whatitmeans{the x-mode sees an opaque wall and the o-mode a thin veil, so most of the light comes out in
one hand of circular polarisation; its sign gives the direction of the field.}
\example{If $\vect B$ points towards the observer, $R$ is the x-mode: $T_R=1.99$\,MK, $T_L=0.36$\,MK, so
$T_b=1.17$\,MK and $P_c=+0.69$. Reversing the field gives $P_c=-0.69$.}
\incode{\code{transfer\_chunk} (\code{coupling="weak"} or \code{"strong"}); \code{gyroresonance/image.py},
class \code{GRImage}: \code{intensity}, \code{stokes\_v}, \code{brightness\_temperature},
\code{tb\_v}, \code{polarization}, \code{flux\_density}; driver \code{rs.make\_gr\_image}.}

%======================================================================
\section{Compact numerical checklist}
%======================================================================
Choose $\nu$ $\rightarrow$ trace one straight ray $\rightarrow$ sample $\vect B$, $n_e$, $T_e$ at the nodes $\rightarrow$
detect $B=B_s$ crossings $\rightarrow$ interpolate the local plasma $\rightarrow$ estimate $L_B$ $\rightarrow$ calculate
$\tau_{s,x}$ and $\tau_{s,o}$ $\rightarrow$ calculate the free--free $\Delta\tau$ $\rightarrow$ assign the modes to $R$ and $L$ $\rightarrow$
perform the radiative transfer (stopping at cutoffs) $\rightarrow$ convert to Stokes $I$ and $V$ $\rightarrow$
repeat for every image pixel.

\begin{center}
\small
\begin{tabular}{@{}l l@{}}
\toprule
Keyword of \code{rs.make\_gr\_image} & Meaning \\
\midrule
\code{mechanisms} (\code{"gr+ff"}, \code{"gr"}, \code{"ff"}) & include layers, free--free, or both \\
\code{s\_max} (default 10) & highest harmonic tested \\
\code{coupling} (\code{"weak"}, \code{"strong"}) & mode coupling, Step 8 \\
\code{gaunt} (\code{"dulk"} or a number) & free--free coefficient, Step 6 \\
\code{ds} & step $\Delta\ell$ along the rays \\
\code{t\_body} & brightness temperature behind the rays, Step 7 \\
\bottomrule
\end{tabular}
\end{center}

\vspace{4pt}
{\small\color{head}\textbf{Source:} \color{gray} Hazra, \emph{Thermal gyroresonance emission in radiosynth:
Equations, derivation and numerical implementation}, radiosynth 2.2.0, November 2026, which gives
the full derivation, the validation tests and the code references. Further reading: Zheleznyakov
1970, \emph{Radio Emission of the Sun and Planets}; Melrose 1980, \emph{Plasma Astrophysics}; White \&
Kundu 1997, Sol.\ Phys.\ 174, 31; Dulk 1985, ARA\&A 23, 169; Cohen 1960, ApJ 131, 664.}

\end{document}
